\slajdid{09_3-s013}
\begin{frame}[label=09_3-s013]
\centering\alert{Jedinični CMOS invertor minimalne geometrije}\par\bigskip\raggedright
U ovom smislu se i definiše \underline{CMOS invertor minimalne geometrije – jedinični CMOS invertor}. Dužine kanala P i N tranzistora su iste, MINIMALNE za datu tehnologiju, a širine zadovoljavaju neke od prethodnih izvedenih odnosa, ili kompromis tih zahteva. Najčešće je odnos 2:1. Zauzima „najmanje prostora“ i kapacitivnosti su minimalne.
\par\medskip\begin{columns}[c,onlytextwidth]\column{.32\textwidth}\centering\input{slike/tikz/s013_jedinicni_invertor.tex}\column{.66\textwidth}
Oznaka 1 i 2, normalizovane širine u datoj tehnologiji.\par\medskip
Ako se tehnologija promeni odnosi ostaju isti.\par\medskip
Zaključci ostaju isti.\par\medskip
Brojne vrednosti se menjaju na primer za otpornosti tranzistora, kašnjenje, itd…, ali je kašnjenje i dalje minimalno, prag odlučivanja na sredini napona napajanja itd…
\end{columns}
\end{frame}
